\subsection{Small-leaf refinements of the eighteen-arrow star scope}

\begin{lemma}[Proved: extending the binary and four-state leaf gates]
\label{CP2:A-star-small-leaf-orbit-extension}
Consider a source star with center $M$ and two nonempty stalks.
Suppose one leaf corner has order at most two. The star cannot
have the rational homology and cohomology ring of $\CP^2$ in
either of the following cases:
\begin{enumerate}
\item $|M|\leq11$ and the other stalk has at most two states,
with no restriction on that other leaf corner;
\item $|M|\leq8$, the other leaf corner has order at most three,
and its stalk has at most four states.
\end{enumerate}
The stalk at the first leaf can have any positive size.
\end{lemma}
\begin{proof}
The rational algebra of a monoid $L$ of order at most two is
semisimple: besides the trivial monoid the possibilities are
$C_2$ and $\{1,e\}$ with $e^2=e$, and in both cases the
algebra is $\mathbb Q\times\mathbb Q$. In particular its
trivial right module is projective. Replace this leaf by a
trivial corner and its left-$L$ stalk $T$ by the orbit set
$L\backslash T$. This is an actual functor, and exact orbit
comparison gives a rational cohomology-ring and homology
isomorphism. More explicitly, the bridge bar
$\mathbb Q\otimes^{\mathbf L}_{\mathbb QL}\mathbb Q[T]
\otimes^{\mathbf L}_{\mathbb QM}\mathbb Q$ is naturally
identified with the ordinary $M$ action bar on $L\backslash T$;
the $L$ corner is rationally acyclic. The other edge and its
intersection $BM$ are unchanged, so the mapping-cylinder
comparison gives the same isomorphism for the whole star.

The orbit stalk remains nonempty and the other stalk does not
change. Apply Proposition~\ref{CP2:A-star-binary-leaf-eleven}
in the first case and
Proposition~\ref{CP2:A-star-three-leaf-four-states} in the
second. Both propositions exclude the rational target, which
is exactly what the orbit comparison transfers.
\end{proof}

\begin{lemma}[Proved: a singleton leaf with a rationally acyclic corner]
\label{CP2:A-star-singleton-small-leaf-retract}
Let a source star have a leaf on a singleton stalk, with
nontrivial leaf monoid $L$ such that $BL$ is rationally acyclic.
Replacing $L$ by the trivial monoid is an actual categorical
retract and a rational homology and cohomology-ring equivalence.
If the original star has type $\CP^2$, so does the retract.
Consequently an eighteen-arrow projective-plane star cannot
have a singleton stalk at a leaf corner of order two, three,
or four.
\end{lemma}
\begin{proof}
All corner actions on the singleton stalk are trivial. Send
every element of $L$ to the identity of the replacement corner,
and keep every other arrow unchanged. This defines a functor;
its section sends the new identity to $1_L$ and fixes all
other arrows. Their composite on the replacement category is
the identity.

Write $E$ for the other full edge. The nerve of the star is
the union of $BE$ and the singleton-edge nerve, meeting in
$BM$, where $M$ is the center. The singleton-edge space is
$BM*BL$, and its replacement is $BM*\{*\}$, a cone. Both
are rationally acyclic: for the first this follows from the
join homology formula and the rational acyclicity of $BL$.
The replacement map is the identity on $BE$ and $BM$.
Mayer--Vietoris therefore makes the actual retraction a
rational homology equivalence in every degree. Its induced
cohomology isomorphism respects products.

Suppose the original space has type $\CP^2$. The categorical
retraction makes the retract's integral homology a direct
summand of the original integral homology, and its fundamental
group a retract of the trivial group. The rational equivalence
forces the degree-two and degree-four integral summands to be
$\mathbb Z$, and all other reduced integral groups vanish.
The split pullbacks in those two degrees are isomorphisms;
naturality consequently transfers the primitive square.
The integral recognition argument in
Theorem~\ref{A:monoid-upper23} gives type $\CP^2$ for the
retract. The replacement removes $|L|-1\geq1$ arrows, so at
eighteen arrows this contradicts
Theorem~\ref{CP2:A-global-lower-eighteen}.
Finally all monoids of orders at most four have rationally
acyclic classifying spaces by Lemma~\ref{CP2:A-five-types}.
\end{proof}

\begin{proposition}[Proved: fifteen further frozen star profiles]
\label{CP2:A-star-eighteen-small-leaf-scope}
In the frozen sixty-one-profile star scope at eighteen arrows,
the original indices
\[
14,17,19,22,24,26,38,40,42,43,44,49,57,58,59
\]
are excluded by the preceding two lemmas. This gives a separate
partition of the $1204$ canonical cardinality profiles into
$1158$ proved profiles and $46$ profiles requiring finite
composition checks. It does not assert that those forty-six
finite checks have all finished.
\end{proposition}
\begin{proof}
The separately frozen program
\path{converse/scope18/star18_small_leaf_supplement.py}
tests each original matrix in both leaf orders. It records
the actual small leaf, the relevant stalk sizes, and the
specific lemma clause. The singleton-corner predicate tests
orders two through four and singleton stalk size. The orbit
predicates test a leaf order at most two and, respectively,
the exact bounds $(|M|,|S|)\leq(11,2)$ or
$(|M|,|N|,|S|)\leq(8,3,4)$ on the retained center, leaf and
stalk. These are precisely the hypotheses proved above.
Exactly the fifteen displayed original rows qualify. The
old ledger and its $1143$ proved rows are unchanged, and
$1143+15=1158$, $61-15=46$.
\end{proof}
